\section{Kisah}

Setiap kali seseorang membuka sebuah halaman web di peramban (\textit{browser}),
terjadi serangkaian proses yang kompleks di balik layar. Peramban mengunduh
dokumen \textit{HyperText Markup Language} (HTML) dari peladen web, kemudian
mengurai dokumen tersebut menjadi sebuah struktur pohon yang disebut
\textit{Document Object Model} (DOM). Struktur pohon inilah yang menjadi
fondasi bagi seluruh proses selanjutnya --- mulai dari penerapan gaya
(\textit{styling}) melalui \textit{Cascading Style Sheets} (CSS), manipulasi
konten secara dinamis melalui JavaScript, hingga proses \textit{rendering}
visual di layar pengguna.

Salah satu mekanisme fundamental dalam pemrosesan halaman web adalah
\textbf{pencocokan CSS Selector}. Ketika seorang pengembang web menuliskan
aturan CSS seperti \texttt{div > p.highlight}, peramban harus menelusuri
seluruh pohon DOM untuk menemukan elemen-elemen yang memenuhi kriteria
tersebut. Proses penelusuran ini pada dasarnya adalah sebuah \textit{tree
traversal} --- mengunjungi node-node dalam pohon secara sistematis hingga
seluruh elemen yang cocok ditemukan.

Dua algoritma traversal yang paling fundamental adalah \textit{Breadth-First
Search} (BFS) dan \textit{Depth-First Search} (DFS). Kedua algoritma ini
memiliki strategi penelusuran yang berbeda: BFS menjelajahi pohon
secara melebar (\textit{level-by-level}), sementara DFS menjelajahi pohon
secara mendalam (\textit{sedalam mungkin baru mundur}). Pemilihan algoritma
traversal yang tepat dapat memengaruhi efisiensi dan urutan penemuan elemen
yang sesuai dengan CSS Selector yang diberikan.

Pada Tugas Besar 2 mata kuliah IF2211 Strategi Algoritma ini, mahasiswa
diminta untuk membangun sebuah aplikasi web yang mengimplementasikan proses
penelusuran CSS Selector pada pohon DOM menggunakan algoritma BFS dan DFS.
Program yang dikembangkan diberi nama \textbf{Berry-Manuka}: aplikasi web
interaktif yang mampu mengambil dokumen HTML dari suatu URL, mem-parse-nya
menjadi pohon DOM, menelusuri pohon tersebut menggunakan BFS maupun DFS untuk
mencari elemen yang cocok dengan CSS Selector tertentu, serta memvisualisasikan
seluruh proses penelusuran secara interaktif.

\section{Latar Belakang}

\textit{Document Object Model} (DOM) adalah antarmuka pemrograman yang
merepresentasikan struktur dokumen HTML dalam bentuk pohon hierarkis.
Setiap elemen HTML dalam dokumen direpresentasikan sebagai sebuah \textit{node}
dalam pohon, dengan hubungan \textit{parent-child-sibling} yang mencerminkan
struktur bersarang (\textit{nesting}) dari tag-tag HTML \cite{mdn_dom}.
Representasi pohon ini bersifat universal dan menjadi dasar bagi hampir semua
operasi yang dilakukan peramban terhadap halaman web.

Pencocokan CSS Selector merupakan operasi yang sangat sering dilakukan oleh
peramban. Setiap kali sebuah halaman web dimuat, peramban harus mencocokkan
setiap aturan CSS dalam \textit{stylesheet} dengan elemen-elemen yang sesuai
di pohon DOM. Proses ini memerlukan penelusuran pohon DOM secara efisien,
mengingat kompleksitas halaman web modern yang dapat memiliki ribuan hingga
puluhan ribu node \cite{mdn_css_selectors}.

Algoritma \textit{Breadth-First Search} (BFS) dan \textit{Depth-First Search}
(DFS) merupakan dua strategi penelusuran graf dan pohon yang paling dasar dan
penting dalam ilmu komputer \cite{cormen2009introduction}. Kedua algoritma ini
memiliki karakteristik yang berbeda dalam hal urutan kunjungan node, penggunaan
memori, dan kesesuaian dengan jenis pencarian tertentu. Pemahaman mendalam
mengenai perbedaan keduanya sangat penting dalam konteks penelusuran pohon DOM.

Visualisasi interaktif dari proses penelusuran pohon DOM dapat membantu
pengguna memahami cara kerja algoritma BFS dan DFS secara lebih intuitif.
Daripada hanya membaca pseudocode atau analisis teoritis, pengguna dapat
mengamati secara langsung bagaimana setiap algoritma mengunjungi node-node
dalam pohon, mengetahui node mana yang diperiksa terlebih dahulu, dan melihat
jalur yang ditempuh hingga elemen yang dicari ditemukan.

Untuk meningkatkan performa penelusuran pada pohon DOM yang berukuran besar,
diterapkan mekanisme konkurensi menggunakan \textit{goroutine} yang disediakan
oleh bahasa pemrograman Go. Dengan pemrosesan paralel, beberapa node dapat
diperiksa secara bersamaan sehingga waktu penelusuran dapat dikurangi secara
signifikan.

Selain itu, diimplementasikan pula algoritma \textit{Lowest Common Ancestor}
(LCA) dengan metode \textit{Binary Lifting} untuk menentukan ancestor bersama
terdekat dari dua buah node dalam pohon DOM. Algoritma ini memiliki
kompleksitas query sebesar $O(\log N)$, jauh lebih efisien dibandingkan
pendekatan naif yang memiliki kompleksitas $O(N)$ per query
\cite{cpalgorithms_lca}.

\section{Deskripsi Tugas}

Program Berry-Manuka merupakan aplikasi web yang menerima masukan berupa
dokumen HTML --- baik melalui \textit{URL scraping} maupun input manual ---
kemudian mem-parse dokumen tersebut menjadi pohon DOM dan menyediakan
fasilitas penelusuran pohon menggunakan algoritma BFS atau DFS berdasarkan
CSS Selector yang diberikan pengguna. Hasil penelusuran divisualisasikan
secara interaktif lengkap dengan animasi proses penelusuran, statistik
performa, dan traversal log.

Program dikembangkan menggunakan bahasa pemrograman Go untuk \textit{backend}
dan framework Next.js untuk \textit{frontend}, serta dikemas dalam
kontainer Docker untuk memudahkan proses deployment. Program ini dikerjakan
secara berkelompok oleh tiga orang mahasiswa.

\subsection{Masukan}

Program menerima beberapa parameter masukan melalui antarmuka web:

\begin{enumerate}
    \item \textbf{Sumber dokumen HTML}: dapat berupa URL website yang akan
    di-\textit{scrape} atau teks HTML yang dimasukkan secara manual oleh
    pengguna.

    \item \textbf{Pilihan algoritma traversal}: pilihan antara
    \textit{Breadth-First Search} (BFS) atau \textit{Depth-First Search}
    (DFS) sebagai metode penelusuran pohon DOM.

    \item \textbf{CSS Selector}: selector yang digunakan untuk mencari elemen
    yang sesuai dalam pohon DOM. Selector yang didukung meliputi
    \textit{tag selector} (misalnya \texttt{p}), \textit{class selector}
    (misalnya \texttt{.highlight}), \textit{ID selector} (misalnya
    \texttt{\#header}), dan \textit{universal selector} (\texttt{*}), serta
    \textit{combinator} berupa child (\texttt{>}), descendant (spasi),
    adjacent sibling (\texttt{+}), dan general sibling (\texttt{\textasciitilde}).

    \item \textbf{Jumlah hasil}: banyaknya kemunculan elemen yang ingin
    ditampilkan, baik berupa $N$ hasil pertama (\textit{top-N}) maupun
    seluruh kemunculan.
\end{enumerate}

\subsection{Keluaran}

Program menghasilkan beberapa bentuk keluaran melalui antarmuka web:

\begin{enumerate}
    \item \textbf{Visualisasi pohon DOM}: representasi visual dari struktur
    pohon DOM hasil \textit{parsing}, lengkap dengan keterangan kedalaman
    maksimum (\textit{maximum depth}) dari pohon tersebut.

    \item \textbf{Highlight jalur traversal}: penandaan visual pada
    node-node yang dikunjungi selama proses penelusuran, sehingga pengguna
    dapat mengamati urutan dan jalur yang ditempuh oleh algoritma.

    \item \textbf{Animasi penelusuran}: animasi langkah demi langkah dari
    proses penelusuran pohon DOM yang menunjukkan node mana yang sedang
    diperiksa pada setiap tahapan.

    \item \textbf{Statistik performa}: informasi mengenai waktu pencarian
    (dalam milidetik) dan banyaknya node yang dikunjungi selama proses
    penelusuran.

    \item \textbf{Traversal log}: catatan berurutan yang memuat informasi
    setiap tahapan penelusuran, termasuk node yang dikunjungi dan status
    pencocokan terhadap CSS Selector.

    \item \textbf{Lowest Common Ancestor}: hasil pencarian ancestor bersama
    terdekat dari dua buah node yang dipilih oleh pengguna dalam pohon DOM,
    beserta jalur dari masing-masing node menuju ancestor tersebut.
\end{enumerate}
